\documentclass[a4paper,10pt]{article}

\usepackage{amsmath}
\usepackage{amssymb}
\usepackage[dvipsnames]{xcolor}
\usepackage{graphicx}
\usepackage{listings}
\usepackage{lmodern}
\usepackage{mdframed}
\usepackage{url}

\usepackage{keyval}

\usepackage{fancyref}
\usepackage{minted}

\newenvironment{code}[3][ht]
 {\VerbatimEnvironment
  \begin{listing}[#1]
  \caption{#2}\label{lst:#3}
  \begin{minted}[
    breaklines,
    escapeinside=||,
    linenos,
    xleftmargin=21pt
    ]{rust}}
 {\end{minted}\end{listing}}

\makeatletter

\define@key{embed}{pos}{\def\e@pos{#1}}
\define@key{embed}{firstline}{\def\e@firstline{#1}}
\define@key{embed}{lastline}{\def\e@lastline{#1}}

\setkeys{embed}{pos=ht,firstline=1,lastline=-1}
  
\newcommand{\embed}[3][]
 {\begingroup
  \setkeys{embed}{#1}
  \expandafter\listing\expandafter[\e@pos]
  \caption{#2}\label{lst:#3}
  \inputminted[
    autogobble,
    breaklines,
    linenos,
    firstline={\e@firstline},
    lastline={\e@lastline},
    firstnumber=1,
    xleftmargin=21pt
    ]{rust}{Code/#3/src/main.rs}
  \endlisting
  \endgroup}

\makeatother

\newcommand*{\fancyreflstlabelprefix}{lst}
\frefformat{vario}{\fancyreflstlabelprefix}{listing\fancyrefdefaultspacing#1#3}
\Frefformat{vario}{\fancyreflstlabelprefix}{Listing\fancyrefdefaultspacing#1#3}

\newmdenv[linecolor=BlueGreen,frametitle=Note,frametitlebackgroundcolor=BlueGreen!40!]{notebox}

\title
{
   \includegraphics[width=12cm]{up_logo.png} \\
   \vspace{2cm}
   \textbf{COS700 Research} \\ \vspace{0.5cm}
   \textbf{Design Pattern Metaprogramming Foundations in Rust\\ \large Abstract Factory and Visitor} \\ \vspace{0.5cm}
   \textbf{Student number:} u19239395 \\ \vspace{0.5cm}
   \textbf{Supervisor}: \\ Dr. Linda Marshall
}

\date{??? 2020}
\begin{document}
\author{}

\maketitle

\newpage
\linespread{1.25}

\section*{Abstract}

\section*{Keywords:}

\newpage

\section{Introduction}

\section{Design Patterns}

\subsection{Abstract Factory}

\subsection{Visitor}

\section{Metaprogramming}
% So far not needed for selected patterns

\section{Rust}
Rust is a relatively new language created by Mozilla to be memory safe yet have low level like performance \cite{klabnik_2019_01}.
Traditionally, memory safe languages will make use of a garbage collector which slows performance \cite{hertz_05_01}.
Garbage collector languages include C\# \cite{robinson_04_01}, Java \cite{gosling_96_01}, Python \cite{martelli_06_01}, Golang \cite{tsoukalos_18_01} and Javascript \cite{flanagan_06_01}.
Languages that perform well use manual memory management, which is not memory safe whenever the programmer is not careful.
Dangling pointers \cite{caballero_12_01}, memory leaks \cite{wilson_92_01}, and double freeing \cite{sharp_13_01} in languages like C and C++ are prime examples of manual memory management problems \cite{konrad_18_01}.
Few languages have both memory safety and performance.
However, Rust achieves both by using a less popular model known as ownership.

\subsection{Ownership}
In the ownership model, the compiler uses statical analysis \cite{rasmussen_2019_01} to track which variable owns a piece of heap data -- this does not apply to stack data.
Each data piece can only be owned by one variable at a time.
The owning variable is called the \textit{owner}.
\cite{klabnik_2019_01}

A variable also has a scope.
The scope starts at the variable declaration and ends at the closing curly bracket of the code block containing the variable.
When the owner goes out of scope, Rust returns the memory by calling the \textit{drop} method at the end of the scope.
Ownership is manifested in two forms -- moving and borrowing.
These two forms are explained next.
\cite{klabnik_2019_01}

\subsubsection{Moving}
Moving happens when one variable is assigned to another.
The compiler's analysis moves ownership of the data to the new variable from the initial variable.
The initial variable's access is then invalidated \cite{klabnik_2019_01}.
An analogy example would be to give a book to a friend.
The friend can do anything from annotating to burning the book as they feel fit since the friend is the book's owner.

\embed[firstline=2,lastline=7]{Example of ownership transfer}{Rust/moved}

In \fref{lst:Rust/moved}, on line 2, a heap data object is created and assigned to variable \textit{s}.
Line 3 assigns \textit{s} to \textit{t}.
However, because \textit{s} is a heap object, the compiler transfers ownership of the data from \textit{s} to \textit{t} and marks \textit{s} as invalid.

When trying to use the data on line 5, via \textit{s}, the compiler, therefore, throws an error saying \textit{s} was moved.
Any reference to \textit{s} after line 3 will always give a compiler error.

Finally, the scope of \textit{t} ends on line 6.
Since the compiler can guarantee \textit{t} is the only variable owning the data, the compiler can free the memory on line 6.

\embed{Function taking ownership}{Rust/fn-move}

Having ownership moving makes excellent memory guarantees within a function but is annoying when calling another function, as seen in \fref{lst:Rust/fn-move}.
The \textit{take\_ownership} function takes ownership of the heap data resulting in memory cleanup code correctly being inserted at the end of \textit{take\_ownership}'s scope on line 10.
When \textit{main} calls \textit{take\_ownership}, \textit{a} becomes the new owner of \textit{s}'s data, making the call on line 5 invalid.
When taking ownership is not desired, the second form of ownership, borrowing, should be used instead.

\subsubsection{Borrowing}

Borrowing has a new variable take references to data rather than becoming its new owner \cite{klabnik_2019_01}.
An analogy is to borrow a book from a friend where a promise is made to return the book to its owner once done with it.

\embed{Function taking borrow}{Rust/fn-borrow}

As seen in \fref{lst:Rust/fn-borrow}, borrowing makes the function \textit{take\_borrow} take a reference to the data.
References are activated with an ampersand (\textit{\&}) before the type.
Once \textit{take\_borrow} has ended, control goes back to \textit{main} - where the cleanup code will be inserted.
Having references as function parameters is called borrowing \cite{klabnik_2019_01}.
The ampersand is also used in the call argument to signal the called function will borrow the data.

However, in Rust, all variables are immutable by default \cite{klabnik_2019_01}.
Hence changing the data in \textit{take\_borrow} causes an error stating the borrow is not mutable on line 9.
Returning to the borrowed book analogy.
One would not make highlights and notes in a book one borrowed unless the owner gave explicit permission.

\subsection{Immutable by default}
Mutable borrows are an explicit indication that a function/variable is allowed to change the data.

\embed{Function taking mutable borrow}{Rust/fn-mut-borrow}

As seen in \fref{lst:Rust/fn-mut-borrow}, mutable borrows are activated using \textit{\&mut } on the type.
Again, \textit{mut} is also used in the call to make it explicit the function will modify the data.
Variables -- on the stack or heap -- also need to be declared \textit{mut} to use them as mutable.
\cite{klabnik_2019_01}

The two ownership forms - moving and borrowing - together with mutable variables put some constraints on the code of each variable type and their calls: \cite{klabnik_2019_01}
\begin{itemize}
  \item Moving will always invalidate the variable.
  \item Borrowed variables cannot be mutated.
        However, more than one function can borrow the data at the same time in parallel and concurrent code.
  \item Mutable borrowing does allow mutations.
        But only one function can hold a mutable borrow at a time and no other immutable borrows can exist.
\end{itemize}

The constraints will always be enforced by the compiler, thus requiring all code to meet them.
Meeting these constraints also requires some shift in thinking at times.
Another shift is required because Rust may not classify as an Object-Oriented Programming (OOP) language.

\subsection{Not quite OOP}
No single definition exists to qualify a language as Object-Oriented \cite{meyer_97_01,stefik_85_01,gamma_94_01,klabnik_2019_01}.
Three Object-Oriented definitions will be explored to understand Rust better.
These three are Objects as Data and their Behaviour, Encapsulation, and Inheritance.

\subsubsection{Objects as Data and their Behaviour}
The first definition of Object-Oriented design is a language using objects.
An object, in turn, holds both data and procedures operating on the data \cite{meyer_97_01,stefik_85_01,gamma_94_01}.
Rust meets the data with their operations requirement by holding the data in \textit{struct}s and having the operations defined in \textit{impl} blocks \cite{klabnik_2019_01}.

\paragraph{Struct}
A \textit{struct} is the same as a \textit{struct} in C \cite{stroustrup_13_01} and other C like languages \cite{robinson_04_01, savitch_15_01, malik_09_01}.
Structs are used to define objects with named data pieces, as shown in \fref{lst:Rust/oop} lines 1 to 5.
Each of the struct properties is named followed by a type.

\embed[lastline=44]{Example of a Struct}{Rust/oop}

\paragraph{Methods}
The operations to perform on a struct are defined in an \textit{impl} block, as seen in lines 7 - 19 in \fref{lst:Rust/oop}.
Notice how the ownership rules apply to the struct.

The methods \textit{get\_age} and \textit{have\_birthday} will take a borrow of the struct object.
To age, while having a birthday, \textit{have\_birthday} needs to take a mutable borrow of \textit{self} - also why \textit{bar} needs to be \textit{mut} in \textit{main}.
The method \textit{have\_burial} moves \textit{self}, thus invalidating any objects of \textit{Foo} after \textit{have\_burial} is called - resulting in a compile error on line 33.
The same error happened in \fref{lst:Rust/fn-move}.

\begin{notebox}
  Rust does not always use the \textit{return} keyword.
  The missing semi-colon on line 9 is deliberate and the Rust way to say we want to return the age for the Foo instance.
\end{notebox}

\subsubsection{Encapsulation}
The next definition deals with hiding implementation details from the client - known as encapsulation \cite{klabnik_2019_01, meyer_97_01}.
Encapsulation allows the struct creator to change the internal procedures of a struct without affecting the public interface used by clients.
Rust also meets the encapsulation definition by using the \textit{pub} keyword.

\paragraph{pub}
As can be seen in \fref{lst:Rust/oop}, the \textit{Foo} struct and its \textit{name} data member is made public explicitly.
The methods \textit{get\_age} and \textit{have\_birthday} are also made public explicitly.
All the other data members and methods are private - Rust's default - and unreachable by client code.
A curious question then is why can \textit{main} access the private variables and methods.
The answer being: \textit{main} is located in the same module/file as the struct and therefore has full access to it.

\begin{notebox}
  Rust does not provide constructors like other OOP languages.
  Instead, Rust has what it calls associate methods.
  An associate method is a method definition not containing \textit{self} in the parameter list \cite{klabnik_2019_01}.
  A constructor like associate-method will return an owned instance of the struct being constructed.
  \fref{lst:Rust/oop}, lines 36 to 44 shows an example of an associate method for Foo defined in the same file as Foo's definition.
  Yes, Rust code can have multiple \textit{impl} blocks for a single struct.
\end{notebox}

\subsubsection{Inheritance}
The last definition being looked at is inheritance.
Inheritance has an object inherit some of its data members and procedures from a parent object \cite{meyer_97_01, stefik_85_01, gamma_94_01}.
Inheritance is mostly used to reduce code duplication.
Rust does not meet the inheritance definition for OOP.

The Gang of Four made their design patterns for Object-Oriented languages.
They also define two principles for Object-Oriented design \cite{gamma_94_01}.
\begin{itemize}
  \item ``Program to an interface, not an implementation.''
        Rust meets this requirement by using \textit{traits}.
  \item ``Favor object composition over class inheritance.''
        The lack of inheritance will, therefore, not be an issue in Rust.
        Not having inheritance causes Rust code to use composition naturally.
\end{itemize}

Rust \textit{trait}s cause it to meet both requirements allowing Rust to implement the Gang of Four design patterns.

% Traits
%% Bounds
%% Defaults

% Sized
% - Because stack is default
%% Opting out

% https://joshleeb.com/blog/rust-traits-trait-objects/
% Box<Trait> is owned; &Trait is a reference.

% macro rules: Hygiene and token list

% Generics
%% Associate types

% Enums
%% Pattern matching

\section{Reporting}

\section{Conclusion}


\bibliographystyle{alpha}
\bibliography{References}

\end{document}